\documentclass[11pt,letterpaper]{article}
\PassOptionsToPackage{hyphens}{url}
\usepackage[margin=.83in,top=.8in,bottom=.8in,headheight=15pt]{geometry}
\usepackage[T1]{fontenc}
\usepackage{lmodern,amsmath,amssymb,mathtools,bm,booktabs,array,graphicx,tikz}
\usetikzlibrary{arrows.meta,positioning,calc,decorations.pathreplacing,patterns}
\definecolor{ink}{HTML}{17344D}\definecolor{accent}{HTML}{087E8B}\definecolor{pale}{HTML}{EDF5F5}\definecolor{gold}{HTML}{AD6B19}
\providecommand{\dd}{\,\mathrm{d}}\providecommand{\R}{\mathbb{R}}
\providecommand{\E}{\mathbb{E}}\providecommand{\Var}{\operatorname{Var}}
\providecommand{\Cov}{\operatorname{Cov}}\providecommand{\diag}{\operatorname{diag}}
\providecommand{\norm}[1]{\left\lVert#1\right\rVert}
\providecommand{\abs}[1]{\left|#1\right|}
\usepackage{microtype,hyperref,fancyhdr,enumitem,amsthm,sectsty,needspace}
\allsectionsfont{\raggedright}
\hypersetup{colorlinks=true,linkcolor=ink,urlcolor=accent,pdfauthor={Zan Ahmad},pdftitle={Repeated observations, group leakage and predictive baselines / Review}}
\setlength{\parindent}{0pt}\setlength{\parskip}{6pt}\linespread{1.07}
\setlist{itemsep=3pt,topsep=4pt,leftmargin=*}
\emergencystretch=2em\clubpenalty=10000\widowpenalty=10000
\pagestyle{fancy}\fancyhf{}\fancyhead[L]{\small Computational Cardiology}\fancyhead[R]{\small Lecture 15}
\fancyfoot[L]{\footnotesize Zan Ahmad}\fancyfoot[R]{\footnotesize\thepage}\renewcommand{\headrulewidth}{.3pt}
\newtheorem{proposition}{Proposition}\newtheorem{theorem}{Theorem}\newtheorem{lemma}{Lemma}
\theoremstyle{definition}\newtheorem{definition}{Definition}\newtheorem{example}{Example}
\numberwithin{equation}{section}
\begin{document}
\section*{Repeated observations, group leakage and predictive baselines --- review sheet}
Use this after attempting the worked examples and practice. Explain each claim in words before calculating.

\section*{Key ideas}
\begin{itemize}
\item Define the target before choosing the split.
\item Audit both rows and groups; exact-row separation is insufficient for unseen people.
\item All fitted preprocessing belongs inside the training boundary.
\item Baselines and class-specific metrics make the claim inspectable.
\end{itemize}
\section*{Common errors and repairs}
\textbf{Treat a perfect contaminated score as new-person performance.} It measures a task with already represented groups; use the target-matched protocol.

\textbf{Fit the scaler on the full table.} Fit on training rows and reuse the stored transform.

\textbf{Treat class balance in this generator as population prevalence.} It is a teaching design choice.

\section*{Retrieve without your notes}
\begin{itemize}
\item What is the proper split's shared-group count?
\item Why does matching training row counts help?
\item Does grouping by person answer every site-shift question?
\end{itemize}
\section*{Readiness check}
\begin{itemize}
\item Define a prediction target and its appropriate unit of separation.
\item Audit both row-index and group-identity overlap in a proposed split.
\item Fit transparent baselines with training-only preprocessing and interpret class-specific metrics.
\item Explain why a strong contaminated score does not estimate performance on new groups.
\end{itemize}
\end{document}
